% ============================================================
%  Axis demo deck — every theme feature in 14 frames
%  Build: pdflatex demo.tex (3 passes)
% ============================================================
\documentclass[aspectratio=169,11pt]{beamer}
\usepackage[utf8]{inputenc}
\usepackage[T1]{fontenc}
\usepackage{lmodern}
\usepackage{amsmath,amssymb,booktabs,graphicx}

\usetheme{axis}

% --- customize (all optional) ---
\title{A Demonstration\\ of the Axis Theme}
\subtitle{Seminar Talk}
\author{A.\ Author}
\date{October 2026}
\axisadvisor{Advisor: Prof.\ A.\ Advisor}
\axiscommittee{Committee: Prof.\ X \quad$\cdot$\quad Prof.\ Y}
% \axislogo{logo.png}          % bottom-right logo (optional)
% \axisset{fa=0.25, fb=0.5, fc=0.75, total=0}

\begin{document}

% ================= TITLE =================
\begin{frame}[plain]
  \titlepage
\end{frame}

\partdivider{Introduction}{Why this theme exists}

% ================= CONTENT =================
\begin{frame}{Lists: Spectral-Stick Bullets}
  \begin{itemize}
    \item Level-one items carry a blue spectral stick.
    \item Wrapped lines re-indent to the text column, and the
          stick never inflates line spacing (it is \emph{smashed}
          to zero height).
      \begin{itemize}
        \item Level two carries a shorter orange stick.
          \begin{itemize}
            \item Level three, a shorter one still.
          \end{itemize}
      \end{itemize}
    \item Enumerated items use bold blue numerals:
      \begin{enumerate}
        \item first,
        \item second.
      \end{enumerate}
  \end{itemize}
\end{frame}

\begin{frame}{Blocks: Flat Tinted Cards}
  \begin{block}{A block}
    Blocks are flat tinted cards --- no beamer title bars.
  \end{block}
  \begin{alertblock}{An alerted block}
    For warnings, problems, and honest caveats.
  \end{alertblock}
  \begin{exampleblock}{An example block}
    For results and takeaways.
  \end{exampleblock}
\end{frame}

\begin{frame}{Equations and Credits}
  The theme sets no math fonts --- Latin Modern throughout:
  \[
    \mathbf{H}(t) \;=\; \mathbf{H}_0 \;+\; \sum_{k} u_k(t)\,\mathbf{C}_k,
    \qquad
    \lvert\psi\rangle \;\mapsto\; e^{-i\mathbf{H}t}\lvert\psi\rangle.
  \]
  \begin{columns}[T,onlytextwidth]
    \begin{column}{0.48\textwidth}
      \begin{itemize}
        \item Two-column layouts align at the top.
        \item Charts, tables, and figures behave as usual.
      \end{itemize}
    \end{column}
    \begin{column}{0.48\textwidth}
      \begin{block}{A block in a column}
        Cards inside columns work too.
      \end{block}
    \end{column}
  \end{columns}
  \credit{Author, Journal \textbf{1}, 123 (2025).}
\end{frame}

\partdivider{Methods}{The pulse waveforms}

\begin{frame}{Divider Waveforms I}
  Divider waveforms are usable anywhere:
  \par\vspace{0.5em}
  \wavegrape\par\vspace{0.6em}
  \wavechirp
  \credit{GRAPE steps; chirp.}
\end{frame}

\begin{frame}{Divider Waveforms II}
  \par\vspace{0.2em}
  \waveramsey\par\vspace{0.5em}
  \waveecho\par\vspace{0.5em}
  \wavespectrum
  \credit{Ramsey pair; echo train; stick spectrum.}
\end{frame}

\partdivider{Results}{Reading the frequency axis}
\partdivider{Outlook}{Fourth part, echo-train waveform}

\begin{frame}{The Frequency-Axis Footline}
  \begin{itemize}
    \item The footline is an axis: hairline with end ticks, part
          boundaries at fractions \texttt{fa}, \texttt{fb}, \texttt{fc}.
    \item Roman numerals label the four parts.
    \item A blue dot marks the current frame:
          $\mathrm{dot} = \mathrm{frame}/\texttt{total}$.
    \item Set the geometry once the deck is final:\\[0.3em]
      \texttt{\textbackslash axisset\{fa=0.236, fb=0.418, fc=0.727, total=55\}}
  \end{itemize}
\end{frame}

\begin{frame}[plain]
  \begin{center}
    \Large\bfseries Thank you\\[0.6em]
    \large Questions?
  \end{center}
\end{frame}

% ================= BACKUP =================
\backupdivider{Backup}{Anticipated questions}

\begin{frame}{B1: A Backup Slide}
  Backup slides keep the axis footline; the divider carries a
  stick spectrum with one highlighted line.
  \credit{Provenance note.}
\end{frame}

\end{document}
